\inicializarSubseccion{Taller de interpolación – Andrés}

\begin{figure}[H]
    \centering
    \begin{minipage}[b]{0.30\textwidth}
        \centering
        \includegraphics[width=\textwidth]{\nat{parte3/andresPuntosUno.png}}
        \caption{Puntos de la interpolación del taller (P1).}
        \label{fig:andresPuntosUno}
    \end{minipage}
    \hfill
    \begin{minipage}[b]{0.30\textwidth}
        \centering
        \includegraphics[width=\textwidth]{\nat{parte3/andresPuntosDos.png}}
        \caption{Puntos de la interpolación del taller (P2) y una función.}
        \label{fig:andresPuntosDos}
    \end{minipage}
\end{figure}

\begin{figure}[H]
    \centering
    \includegraphics[width=0.30\textwidth]{\nat{parte3/andresFunciones.png}}
    \caption{Los polinomios generados por la interpolación.}
    \label{fig:andresFunciones}
\end{figure}

\begin{figure}[H]
    \centering
    \includegraphics[width=0.45\textwidth]{\nat{parte3/andresPlanoGeogebra.png}}
    \caption{Gráfica generada por los puntos en Geogebra.}
    \label{fig:andresPlanoGeogebra}
\end{figure}

\begin{figure}[H]
    \centering
    \includegraphics[width=0.45\textwidth]{\nat{parte3/andresPlanoMATLAB.png}}
    \caption{Gráfica generada por los puntos en MATLAB.}
    \label{fig:andresPlanoMATLAB}
\end{figure}